% --- LaTeX Homework Template - S. Venkatraman ---

% --- Set document class and font size ---

\documentclass[letterpaper, 11pt]{article}

% --- Package imports ---

\usepackage{
  amsmath, amsthm, amssymb, mathtools, dsfont,	  % Math typesetting
  graphicx, wrapfig, subfig, float,                  % Figures and graphics formatting
  listings, color, inconsolata, pythonhighlight,     % Code formatting
  fancyhdr, sectsty, hyperref, enumerate, enumitem } % Headers/footers, section fonts, links, lists

% --- Page layout settings ---

% Set page margins
\usepackage[left=0.7in, right=0.7in, bottom=1in, top=1.1in, headsep=0.2in]{geometry}

% Anchor footnotes to the bottom of the page
\usepackage[bottom]{footmisc}



% Set line spacing
\renewcommand{\baselinestretch}{1.2}

% Set spacing between paragraphs
\setlength{\parskip}{1.5mm}

% Allow multi-line equations to break onto the next page
\allowdisplaybreaks

% Enumerated lists: make numbers flush left, with parentheses around them
\setlist[enumerate]{wide=0pt, leftmargin=21pt, labelwidth=0pt, align=left}
\setenumerate[1]{label={(\arabic*)}}

% --- Page formatting settings ---

% Set link colors for labeled items (blue) and citations (red)
\hypersetup{colorlinks=true, linkcolor=blue, citecolor=red}

% Make reference section title font smaller
\renewcommand{\refname}{\large\bf{References}}

% --- Settings for printing computer code ---

% Define colors for green text (comments), grey text (line numbers),
% and green frame around code
\definecolor{greenText}{rgb}{0.5, 0.7, 0.5}
\definecolor{greyText}{rgb}{0.5, 0.5, 0.5}
\definecolor{codeFrame}{rgb}{0.5, 0.7, 0.5}

% Define code settings
\lstdefinestyle{code} {
  frame=single, rulecolor=\color{codeFrame},            % Include a green frame around the code
  numbers=left,                                         % Include line numbers
  numbersep=8pt,                                        % Add space between line numbers and frame
  numberstyle=\tiny\color{greyText},                    % Line number font size (tiny) and color (grey)
  commentstyle=\color{greenText},                       % Put comments in green text
  basicstyle=\linespread{1.1}\ttfamily\footnotesize,    % Set code line spacing
  keywordstyle=\ttfamily\footnotesize,                  % No special formatting for keywords
  showstringspaces=false,                               % No marks for spaces
  xleftmargin=1.95em,                                   % Align code frame with main text
  framexleftmargin=1.6em,                               % Extend frame left margin to include line numbers
  breaklines=true,                                      % Wrap long lines of code
  postbreak=\mbox{\textcolor{greenText}{$\hookrightarrow$}\space} % Mark wrapped lines with an arrow
}

% Set all code listings to be styled with the above settings
\lstset{style=code}

% --- Math/Statistics commands ---

% Add a reference number to a single line of a multi-line equation
% Usage: "\numberthis\label{labelNameHere}" in an align or gather environment
\newcommand\numberthis{\addtocounter{equation}{1}\tag{\theequation}}

% Shortcut for bold text in math mode, e.g. $\b{X}$
\let\b\mathbf

% Shortcut for bold Greek letters, e.g. $\bg{\beta}$
\let\bg\boldsymbol

% Shortcut for calligraphic script, e.g. %\mc{M}$
\let\mc\mathcal

% \mathscr{(letter here)} is sometimes used to denote vector spaces
\usepackage[mathscr]{euscript}

% Convergence: right arrow with optional text on top
% E.g. $\converge[w]$ for weak convergence
\newcommand{\converge}[1][]{\xrightarrow{#1}}

% Normal distribution: arguments are the mean and variance
% E.g. $\normal{\mu}{\sigma}$
\newcommand{\normal}[2]{\mathcal{N}\left(#1,#2\right)}

% Uniform distribution: arguments are the left and right endpoints
% E.g. $\unif{0}{1}$
\newcommand{\unif}[2]{\text{Uniform}(#1,#2)}

% Independent and identically distributed random variables
% E.g. $ X_1,...,X_n \iid \normal{0}{1}$
\newcommand{\iid}{\stackrel{\smash{\text{iid}}}{\sim}}

% Equality: equals sign with optional text on top
% E.g. $X \equals[d] Y$ for equality in distribution
\newcommand{\equals}[1][]{\stackrel{\smash{#1}}{=}}

% Math mode symbols for common sets and spaces. Example usage: $\R$
\newcommand{\R}{\mathbb{R}}   % Real numbers
\newcommand{\C}{\mathbb{C}}   % Complex numbers
\newcommand{\Q}{\mathbb{Q}}   % Rational numbers
\newcommand{\Z}{\mathbb{Z}}   % Integers
\newcommand{\N}{\mathbb{N}}   % Natural numbers
\newcommand{\F}{\mathcal{F}}  % Calligraphic F for a sigma algebra
\newcommand{\El}{L} % Calligraphic L, e.g. for L^p spaces

% Math mode symbols for probability
\newcommand{\pr}{\mathbb{P}}    % Probability measure
\newcommand{\E}{\mathbb{E}}     % Expectation, e.g. $\E(X)$
\newcommand{\var}{\text{Var}}   % Variance, e.g. $\var(X)$
\newcommand{\cov}{\text{Cov}}   % Covariance, e.g. $\cov(X,Y)$
\newcommand{\corr}{\text{Corr}} % Correlation, e.g. $\corr(X,Y)$
\newcommand{\B}{\mathcal{B}}    % Borel sigma-algebra

% Other miscellaneous symbols
\newcommand{\tth}{\text{th}}	% Non-italicized 'th', e.g. $n^\tth$
\newcommand{\Oh}{\mathcal{O}}	% Big-O notation, e.g. $\O(n)$
\newcommand{\1}{\mathds{1}}	% Indicator function, e.g. $\1_A$

% Additional commands for math mode
\DeclareMathOperator*{\argmax}{argmax}    % Argmax, e.g. $\argmax_{x\in[0,1]} f(x)$
\DeclareMathOperator*{\argmin}{argmin}    % Argmin, e.g. $\argmin_{x\in[0,1]} f(x)$
\DeclareMathOperator*{\spann}{Span}       % Span, e.g. $\spann\{X_1,...,X_n\}$
\DeclareMathOperator*{\bias}{Bias}        % Bias, e.g. $\bias(\hat\theta)$
\DeclareMathOperator*{\ran}{ran}          % Range of an operator, e.g. $\ran(T) 
\DeclareMathOperator*{\dv}{d\!}           % Non-italicized 'with respect to', e.g. $\int f(x) \dv x$
\DeclareMathOperator*{\diag}{diag}        % Diagonal of a matrix, e.g. $\diag(M)$
\DeclareMathOperator*{\trace}{trace}      % Trace of a matrix, e.g. $\trace(M)$

% Numbered theorem, lemma, etc. settings - e.g., a definition, lemma, and theorem appearing in that 
% order in Section 2 will be numbered Definition 2.1, Lemma 2.2, Theorem 2.3. 
% Example usage: \begin{theorem}[Name of theorem] Theorem statement \end{theorem}
\theoremstyle{definition}
\newtheorem{theorem}{Theorem}[section]
\newtheorem{proposition}[theorem]{Proposition}
\newtheorem{lemma}[theorem]{Lemma}
\newtheorem{corollary}[theorem]{Corollary}
\newtheorem{definition}[theorem]{Definition}
\newtheorem{example}[theorem]{Example}
\newtheorem{remark}[theorem]{Remark}

% Un-numbered theorem, lemma, etc. settings
% Example usage: \begin{lemma*}[Name of lemma] Lemma statement \end{lemma*}
\newtheorem*{theorem*}{Theorem}
\newtheorem*{proposition*}{Proposition}
\newtheorem*{lemma*}{Lemma}
\newtheorem*{corollary*}{Corollary}
\newtheorem*{definition*}{Definition}
\newtheorem*{example*}{Example}
\newtheorem*{remark*}{Remark}
\newtheorem*{claim}{Claim}
\newcommand{\Tau}{\mathcal{T}}
\newcommand{\PSet}{\mathcal{P}}


% --- Left/right header text (to appear on every page) ---

% Include a line underneath the header, no footer line
\pagestyle{fancy}
\renewcommand{\footrulewidth}{0pt}
\renewcommand{\headrulewidth}{0.4pt}

% Left header text: course name/assignment number
\lhead{MAM4000W, Analysis II -- Assignment 2}

% Right header text: your name
\rhead{Alex Cristaudo, CRSALE010}

% --- Document starts here ---

\begin{document}
\subsection*{Question 1}

\begin{enumerate}
	\item [(a)] We have that \[\{f_n > t\} = \left \{x \in \mathbb{R} : f \left (\frac{[nx]}{n} \right ) > t \right \} = \left \{x \in \mathbb{R} : \frac{[nx]}{n} \in \{f>t\} \right \} = \{x \in \mathbb{R} : [nx] \in n\{f > t\}\} \]\[= \{x \in \mathbb{R}: \exists k \in n\{f>t\}, [nx]=k\} \]
	But for any $k \in n\{f>t\}$, we have that $[nx] = k \iff k \in \mathbb{Z} \text{ and } k \le nx < k+1 \iff k \in \mathbb{Z} \text{ and } \\ x \in \left [ \frac{k}{n}, \frac{k+1}{n}\right )$. So
	\[
	\{f_n > t\} = \bigcup \limits _{\substack{k \in n\{f>t\} \\ k \in \mathbb{Z}}} \left [ \frac{k}{n}, \frac{k+1}{n}\right )
	\]
	
	\item [(b)] Each $\left [ \frac{k}{n}, \frac{k+1}{n}\right ) \in \mathcal{B}(\mathbb{R})$, so $\bigcup \limits 
	_{\substack{k \in n\{f>t\} \\ k \in \mathbb{Z}}} \left [ \frac{k}{n}, \frac{k+1}{n}\right ) \in \mathcal{B}(\mathbb{R})$, as a countable union so $f_n$ is Lebesgue measurable
	
	\item [(c)] We have that $\frac{[nx]}{n} \le \frac{nx}{n} = x$. So the sequence $\left (\frac{[nx]}{n} \right )_{n \in \mathbb{N}}$ is an increasing sequence converging to $x$. So, since $f$ is continuous from the left, we have that 
	\[
	\lim \limits _{n \to \infty} f_n(x) = \lim \limits _{n \to \infty} f \left ( \frac{[nx]}{n}\right ) = \lim \limits _{y \to x^-} f(y) = f(x)
	\]
	Then, since each $f_n$ is Lebesgue measurable, $f$ is as a limit of Lebesgue measurable functions.
\end{enumerate}

\subsection*{Question 2}
\begin{enumerate}
	\item [(a)] We have that \[
	\int _{(1, \infty)} |f|^p d\lambda _1 = \int _{\mathbb{R}} |f|^p \chi _{(1, \infty)} d \lambda _1 = \int _{\mathbb{R}} \lim \limits _{n\to \infty} (\chi _{(1, n)} |f|^p) d \lambda _1
	\]
	The function $|f|^p \chi _{(1, n)}$ is increasing, positive, and measurable ($|f|^pL$ is measurable since it is continuous on $(1, n)$, $\chi _{(1,n)}$ is measurable since $(1,n)$ is a Borel set, so $|f|^p \chi _{(1,n)}$ is measurable as a product of measurable function) so by the Monotone Convergence Theorem:
	\[
	\int _{(1, \infty )} |f|^p d \lambda _1 = \lim \limits _{n \to \infty} \int_{\mathbb{R}} |f|^p \chi _{(1,n)} d\lambda _1 = \underbrace{\lim _{n \to \infty} \int _{(1,n)} |f|^p d \lambda _1 = \lim \limits _{n \to \infty} \int _1^n |f|^p dx}_{f \text{ is bounded on } (1,n) \text{ so Lebesgue = Riemann }} = \int _{1} ^\infty |f|^p dx
	\]
	So we can consider the improper Riemann integral and use the Comparison Test. Since $1+ |\ln (x)| \ge 1$, we get \[\frac{1}{(1+|\ln (x)|)^p x^{p/2}} \le \frac{1}{x^{p/2}}\]Now \[\int _{1}^\infty \frac{1}{x^{p/2}} dx \text{ converges} \iff \frac{p}{2} > 1 \iff p > 2\]
	So, for $p > 2$, we get that $f \in L^p(1, \infty)$. On the other hand, we can write $(1, \infty) = \bigcup \limits _{n \in \mathbb{N}^*} (e^n, e^{n+1}]$. So, if we can show \[
	 \lim \limits _{n \to \infty} \int  _1 ^{e^{n+1}} \frac{1}{(1+|\ln (x)|)^p x^{p/2}} dx= \lim \limits _{n \to \infty} \sum \limits _{k=0}^{n} \int _{e^k}^{e^{k+1}} \frac{1}{(1+|\ln (x)|)^p x^{p/2}}dx
	\]diverges, then so will $\int _1 ^ \infty \frac{1}{(1+|\ln (x)|)^p x^{p/2}} dx$. On each interval $(e^k, e^{k+1}], \ln x \le k+1$ and $x \le e^{k+1}$. So,
	\[
	\sum \limits _{k=0}^{n} \int _{e^k}^{e^{k+1}} \frac{1}{(1+|\ln (x)|)^p x^{p/2}}dx \ge \sum \limits _{k=0}^{n} \int _{e^k}^{e^{k+1}} \frac{1}{(k+2)^p e^{\frac{p(k+1)}{2}}}dx = \sum \limits _{k=0}^{n} \frac{e^k(e-1)}{(k+2)^p e^{\frac{p(k+1)}{2}}} \ge \sum \limits _{k=0}^n \frac{1}{(k+2)^pe^{\frac{1}{2}(pk + p - 2k)}}
		\]
	If $p < 2$, then $pk + p - 2k = k(p - 2) + p$, which, for $k$ large enough, is negative. Then, 
	\[	\lim _{k \to \infty} \frac{1}{(k+2)^pe^{\frac{1}{2}(pk + p - 2k)}} = \infty
	\] so this series diverges, meaning $f \notin L^p (1, \infty)$

	\item [(b)] We have that \[
	\int \limits _{1} ^\infty \frac{1}{(1+\ln x)^2 x} dx  = \lim \limits _{n \to \infty} \int \limits _1 ^n \frac{1}{(1+\ln x)^2 x}dx
	\]
	Letting $u = \ln x$, we get \[
	\int \limits _{1} ^n \frac{1}{(1+\ln x)^2 x} dx  = \int _0 ^{\ln n} \frac{1}{(1+u)^2}dx = - \frac{1}{1+u} \bigg ]_0^ {\ln n} = 1-\frac{1}{1+ \ln n} \to 1 \text{ as } n\to \infty
	\]
	So $f \in L^2(1, \infty)$
	
	\item[(c)] Again, with the same reasoning as in (a), let $(a_n)$ be any decreasing sequence converging to $0$. Then we have \[
	\int _{(0,1)} |f|^p d \lambda _1 = \int _\mathbb{R} |f|^p \chi _{(0,1)} d \lambda _1 = \int _{\mathbb{R}} \lim \limits _{n \to \infty} (|f|^p \chi _{(a_n, 1)}) d\lambda _1 = \lim \limits _{n \to \infty} \int _{(a_n, 1)} |f|^p d \lambda _1 = \lim \limits _{n \to \infty} \int _{a_n} ^1 |f|^p dx
	\]
	Since this is true for any sequence $(a_n)$, we get \[
	\int _{(0,1)} |f|^p d \lambda _1 = \int _0 ^1 |f|^p dx
	\]
	Again, we can consider the improper Riemann integral and use comparison tests.
	
Again,  $\frac{1}{(1+|\ln (x)|)^p x^{p/2}} \le \frac{1}{x^{p/2}}$. Now \[
	\int _0 ^1 \frac{1}{x^{p/2}}dx = \lim \limits _{a \to 0^+} \int \limits _a^1 x^{-p/2}dx = \begin{cases} \ln (x) \bigg ] _a ^1 & \text{ if } p/2 = 1 \\ \frac{1}{1-p/2} x^{1-p/2} \bigg ] _a ^1 & \text{ if } p/2 \ne 1\end{cases}
	\]
	Then $\ln x \to -\infty$ as $a \to -\infty$. On the other hand, if $1-p/2 > 0$, then $x^{1-p/2} $ converges as $x \to 0^+$ (and diverges otherwise). That is, $f \in L^p(0,1)$ if $p<2$. For $p=2$, we get 
	 \[
	\int \limits _{0} ^1 \frac{1}{(1+\ln x)^2 x} dx  = \lim \limits _{a \to 0^+} \int \limits _a ^1 \frac{1}{(1+\ln x)^2 x}dx
	\]
	Letting $u = \ln x$, we get \[
	= \int _{\ln a} ^{1} \frac{1}{(1+u)^2}dx = - \frac{1}{1+u} \bigg ]_a^ {1} = \frac{1}{1+a}- \frac{1}{2} \to \frac{1}{2} \text{ as } a\to 0^+
	\]
	So $f \in L^2(0, 1)$
	
	On $(0,1)$, for $p > 2$ we have that \[
	\frac{1}{(1+|\ln x|)^p x^{p/2}} \ge \frac{1}{(1+|\ln x|)^p} \ge \frac{1}{1+| \ln x|} = \frac{1}{1-\ln x}. 
	\]
	Noticing that $(0,1) = \bigcup \limits _{n \in \mathbb{N}^*}(1-e^{-n}, 1-e^{-n-1}]$ and doing the same as in (b), we get
	\[
	 \lim \limits _{n \to \infty} \sum \limits _{k=0}^n \int _{1-e^{-k}}^{1-e^{-k-1}} \frac{1}{1-\ln x}dx \ge  \lim \limits _{n \to \infty} \sum \limits _{k=0}^n \int _{1-e^{-k}}^{1-e^{-k-1}} e^kdx = \lim \limits _{n \to \infty} \sum \limits _{k=0}^n e^{-k-1}(e-1)e^k
	\]
	And the terms in the sum $e^{-k-1}e^k(e-1) = 1-e^{-1} \not \to 0$ so this sum diverges. Thus $f \notin L^p(0,1)$ for $p > 2$
	
	\item [(d)] $f \in L^p(0,\infty) \iff \int _0 ^ \infty |f|^p dx < \infty$ . Since we can write $\int _0^ \infty |f|^p dx = \int _0 ^1 |f|^p dx + \int _1 ^\infty |d|^pdx$, it follows $f \in L^p(0, \infty) \iff f \in L^p(0,1) \cap L^p(1, \infty)$. Using our results from (a), (b) and (c), $f \in L^p(0, \infty) \iff p = 2$
\end{enumerate}
\end{document}
